\documentclass[10pt,a4paper]{amsart}
\usepackage[margin=19mm]{geometry}
\usepackage{amsmath,amssymb,mathtools}
\usepackage[scr=rsfs,cal=euler]{mathalfa}
\usepackage{xfrac}
\usepackage[unicode,hidelinks,bookmarks=false]{hyperref}
\newcommand{\ie}{i.e.\ }
\newcommand{\cf}{cf.\ }
\newcommand{\resp}{respectively\ }
\newcommand{\Subsection}{\subsection}
\providecommand{\nobreakdash}{\nobreak\hbox{-}}
\setlength{\emergencystretch}{2em}
\multlinegap0pt
\usepackage{amssymb}
\usepackage{mathtools}
\usepackage{enumerate}
\newtheorem{lemma}{Lemma}
\newcommand{\R}{\mathbb{R}}
\renewcommand{\ge}{\geqslant}
\title{Stability and Self-Organized Patterns \\ in Coupled Ecohydrological--Fire Dynamics: \\ A Model of Vegetation--Rainfall--Bushfire Interactions}
\author{Serena Dipierro, Enrico Valdinoci}
\date{Source: arXiv:2509.04766v2 (CC BY 4.0)}
\begin{document}
\maketitle

\noindent\emph{Source and licence.} Stability and Self-Organized Patterns \\ in Coupled Ecohydrological--Fire Dynamics: \\ A Model of Vegetation--Rainfall--Bushfire Interactions. Version arXiv:2509.04766v2 (CC BY 4.0), distributed by arXiv under the Creative Commons Attribution 4.0 International licence, http://creativecommons.org/licenses/by/4.0/, as recorded in the arXiv licence metadata for this exact version and shown on the abstract page https://arxiv.org/abs/2509.04766v2. Full licence notice: \texttt{licenses/arxiv-2509.04766-CC-BY-4.0.txt}.\par\smallskip

\noindent\textit{Editorial scope.} Counted unit: the article's lemma LE:A (source0.tex lines 932--934). The article's complete original proof of it is delivered with it (source0.tex lines 936--992, one \texttt{proof} environment of 57 lines). No mathematics is written, repaired or re-derived: the statement and the proof are the article's own lines, in the article's own notation, and no retained line is substituted or re-flowed.  No further source-local context is needed: every cross-reference of every retained line resolves inside the two delivered spans.  Nothing was omitted from any retained span, so the package carries no omission record.  No retained line contains a citation, so the wrapper carries no bibliography and no bibliography entry.  No retained line contains a figure environment, an image-inclusion command or drawing code, so no image is delivered and the package declares no asset dependency.  Editorial formatting only, in addition: an AMS \texttt{amsart} preamble that restates the article's own declarations and macros used by the retained lines, namely the environments \texttt{lemma} on separate counters, together with the standard packages those lines need.  The retained environments print in this document as Lemma 1 in order of appearance, whereas the article prints the same items with the numbering of its own full document.  The complete native source of the article as distributed in the arXiv e-print for this exact version is bundled unchanged as \texttt{sources/184506/source0.tex}.
\begin{lemma}\label{LE:A}
Let~$a_0$, $a_1$, $a_2\in(0,+\infty)$ and consider the polynomial $$P(t) = t^3 + a_2 t^2 + a_1 t + a_0.$$ Then, all the roots of~$P$ have negative real part if and only if $a_1 a_2 > a_0 $.
\end{lemma}

\begin{proof}
This is in fact a particular case of the Routh-Hurwitz Stability Criterion, but we provide
here a self-contained proof for the facility of the reader.

To prove Lemma~\ref{LE:A}, we consider the complex roots~$r_1$, $r_2$, and~$r_3$ of~$P$
and the corresponding factorization
$$P(t)=(t-r_1)(t-r_2)(t-r_3),$$
leading to
\begin{equation*}
\begin{split}&
a_2=-(r_1+  r_2 + r_3) ,\\&
a_1=r_1 r_2 + r_1 r_3 + r_2 r_3 , \\ {\mbox{and }}\quad&a_0=-r_1 r_2 r_3 .
\end{split}\end{equation*}

As a result,
\begin{equation}\label{mag}
a_0-a_1a_2=(r_1 + r_2) (r_1 + r_3) (r_2 + r_3).
\end{equation}

We also remark that~$P(+\infty)=+\infty$ and~$P(-\infty)=-\infty$, therefore
\begin{equation}\label{ser.pqkdoewvnb9m87}
{\mbox{$P$ has at least one real root.}}\end{equation} For this reason, without loss of generality, we can suppose that~$r_1$ is real.

Also, since~$P(t)\ge a_0>0$ for all~$t\ge0$, we see that
\begin{equation}\label{ser}
{\mbox{all real roots must be negative}}\end{equation} and, in particular, $r_1<0$.

Now, suppose that all the roots are real and negative,
namely~$r_1$, $r_2$, $r_3\in(-\infty,0)$.
Then, it holds that~$r_i+r_j<0$ for all~$i$, $j\in\{1,2,3\}$ and thus
we infer from~\eqref{mag} that
\begin{equation}\label{Swq90nc79m0bv37}
a_0-a_1a_2<0.\end{equation}

Let us now suppose that~$P$ has complex roots.
In this case, the roots must be conjugated, since the coefficients of~$P$ are real,
hence, in this case, one can write~$r_2=x+iy$ and~$r_3=x-iy$, with~$x$, $y\in\R$,
and we obtain from~\eqref{mag} that, as long as~$x\ne0$,
\begin{equation}\label{ser2} \frac{a_0-a_1a_2}{2x}=(r_1 + x+iy) (r_1 + x-iy)=
2 r_1 x + r_1^2 + x^2 + y^2.\end{equation}

Hence, if~$P$ possesses complex roots with negative real part, we deduce from~\eqref{ser2} that
$$ \frac{a_0-a_1a_2}{2x}\ge 2 r_1 x>0$$
and therefore~$a_0-a_1a_2<0$.

These observations show that if~$P$ has complex
roots with negative real part, then necessarily~$a_1 a_2 > a_0 $.

Suppose now that~$a_1 a_2 > a_0 $. Our goal is to show that all the roots have negative
real part. The claim is true if all the roots are real, due to~\eqref{ser}, therefore
we can suppose that~$P$ has complex roots. In this situation, by~\eqref{ser2},
$$\frac{a_0-a_1a_2}{2x}=
(r_1 + x)^2 + y^2\ge0
$$and therefore~$x<0$, as desired.

The proof of Lemma~\ref{LE:A} is thereby complete.
\end{proof}

\end{document}
